\subsection{Cofinal higher quotients and the first further coefficient}

All commutator subgroups in this subsection are closed, and all group
quotients are pro-three. Put
\[
 R_{\rm met}=N'',\quad R_{\rm mark}=\gamma_3(I_{\rm mark}),\quad
 Q_R=N/R,\quad Q_{R,r}=N/(R\gamma_rN)\quad(r\ge3).
\]
We also use $R=1$ for the full group. The notation
$\operatorname{Sec}_S$ means cocycles inflated from
$G_{\mathbb Q,S}$, $S=\{3,\infty\}$, with its actual action.
The manuscript comparison is \emph{La Longue Marche}, \S51.16,
folios 733v--736v, scans 172--178, formulas (371)--(379).

\begin{statement}{A.20}{Proved}{Every unramified class-two section lifts to the full group}
For each of the three choices of $R$, the map
\[
 \operatorname{Sec}_S(Q_R)\longrightarrow
 \operatorname{Sec}_S(N/\gamma_3N)
\]
is surjective. Its fibers have the following exact description.
Use $u,c$ of A.10 and let $B$ be a pro-three quotient of $N$.
Writing $B^c$ for the fixed subgroup of genuine complex conjugation,
we have a bijection
\[
 \operatorname{Sec}_S(B)=
 B/\{x\sim h^{-1}x\,u(h):h\in B^c\}.                 \tag{A.20.1}
\]
For $p:B\twoheadrightarrow N/\gamma_3N$, put $J=\ker p$.
The fiber over the class represented by $x$ in (A.20.1) is
\[
 p^{-1}(x)/\{y\sim h^{-1}y\,u(h):h\in J^c\}.        \tag{A.20.2}
\]
These orbit sets specify all identifications between the lifts.
\end{statement}
\begin{proof}
The based action on $N$ is unramified outside three. Over $K_0$
its image on $N/\Phi(N)=V/3$ is a three-group: the toric and
elliptic graded pieces are trivial and the extension is unipotent.
The kernel of $\operatorname{Aut}(N)\to\operatorname{Aut}(V/3)$
is pro-three. Indeed its image on a characteristic finite
three-group quotient acts trivially on the Frattini quotient,
and successive lower exponent-three factors show that this
automorphism group has three-power order. Thus the whole action
over $K_0$ factors through $P_3$ of A.10. For a finite
three-group-valued cocycle, its graph together with this action
also has three-group image. The cocycle therefore factors through
the same relative quotient $\Gamma=P_3\rtimes\langle c\rangle$.
Taking characteristic inverse limits proves this for $B$.

We have $H^1(C_2,B)=1$. For finite three-groups, conjugacy of
complements of order two follows inductively by averaging in a
central elementary abelian subgroup. For $B$, the solution sets
are nonempty closed subsets of the compact group, so the finite
intersection property gives the same result. Normalize a
cocycle by $z(c)=1$. Since $u,cuc$ freely generate $P_3$, the
value $z(u)$ is arbitrary and determines $z$ uniquely.
The conjugations preserving this normalization are precisely
$B^c$, and their action is (A.20.1).

For a surjection with this involution, $B^c\to(B/J)^c$ is
surjective: the defect of any lift of a fixed element is a
$C_2$-cocycle in $J$, which the same argument removes.
Lifting $x$ now proves surjectivity on section classes, and
the fixed-subgroup assertion moves every representative of
the fiber above $x$ itself. Its remaining conjugators project
to the stabilizer of $x$ under $(N/\gamma_3N)^c$.
That stabilizer centralizes the whole $\Gamma$-section; it is
trivial by C.1 applied to the inflated arithmetic section.
Thus exactly $J^c$ remains, giving (A.20.2).
\end{proof}

\begin{statement}{A.21}{Proved}{Cofinality and the exact comparison of section limits}
For every number field $L$ and each specified $R$, the natural
maps are an isomorphism of pro-three groups and a bijection of
pointed sets, respectively:
\[
 Q_R\xrightarrow{\sim}\varprojlim_{r\ge3}Q_{R,r},\qquad
 H^1_{\mathrm{cts}}(G_L,Q_R)\xrightarrow{\sim}
 \varprojlim_{r\ge3}H^1_{\mathrm{cts}}(G_L,Q_{R,r}).    \tag{A.21.1}
\]
Each $Q_{R,r}$ has a finite central series of finitely generated
$\mathbb Z_3$-module factors, and every finite quotient of $Q_R$
factors through one of them. They are cofinal finite-rank
quotients for this section problem.
\end{statement}
\begin{proof}
A finite three-group is nilpotent: the class equation gives a
nontrivial center, and induction gives a central series.
Consequently every finite quotient kills some $\gamma_rQ_R$.
This proves cofinality and that their intersection is trivial.
Compactness gives the first isomorphism; a finite collection
of compatible coordinates lifts to $Q_R$, proving surjectivity.
The lower-central factors of $N$ are generated over $\mathbb Z_3$
by the finitely many iterated commutators of fixed degree in its
ten generators. Their quotients give the finite-rank assertion.

For compatible cohomology classes, choose representatives
successively. The image of the next cocycle is conjugate to
the preceding one; lift the conjugator through the surjective
group map and modify the new cocycle. This gives exactly
compatible continuous representatives. Their inverse limit
is continuous for the inverse limit topology and is a cocycle.
For injectivity, the sets of elements of $Q_R$ conjugating two
cocycles at level $r$ are nonempty closed nested subsets of a
compact group. Their intersection supplies one conjugator at
every level. This proves the second assertion of (A.21.1).
\end{proof}

\begin{statement}{A.22}{Proved}{Metabelian coefficient lattices and the next lifting fiber}
For $r\ge2$ put $M_r=\gamma_r(N/N'')/\gamma_{r+1}(N/N'')$.
There is a natural $G_{\mathbb Q}$-equivariant isomorphism
\[
 M_r\simeq\ker\bigl(V\otimes\operatorname{Sym}^{r-1}V
             \longrightarrow\operatorname{Sym}^rV\bigr),\qquad
 \operatorname{rank}M_r=(r-1)\binom{r+8}{r},             \tag{A.22.1}
\]
where the arrow is multiplication. For $r\ge3$, write
$A_{r,q}=M_r(-1)$ and $D_{r,q}=\operatorname{Frob}_q-1$.
Then
\[
\begin{split}
 H^2_{\mathrm{cts}}(G_{\mathbb Q},M_r)
   &=\mathcal R(A_r)/D_r\mathcal R(A_r),\\
 \varprojlim{}^1H^1(G_{\mathbb Q},M_r/3^n)
   &=\frac{\mathcal R(A_r)\cap\prod_{q\ne3}D_{r,q}A_{r,q}}
           {D_r\mathcal R(A_r)}\ne0,\\
 H^1_{\mathrm{cts}}(G_{\mathbb Q},M_r)
   &=M_r/(u-1)M_r^c.
\end{split}
\tag{A.22.2}
\]
The matrices on $M_r$ are the kernel functor (A.22.1) applied
to the explicit matrices on $V$ in A.7, giving a finite integral
prescription for every first-cohomology Smith factor.

The next coefficient of the full, metabelian, and marked
depth-two towers is the same rank-$330$ lattice $M_3$, and
\[
\begin{split}
 H^1(G_{\mathbb Q},M_3)\simeq\mathbb Z_3^{166}
 &\oplus(\mathbb Z/3)^7\oplus(\mathbb Z/9)^{56}\\
 &\oplus(\mathbb Z/27)^{25}\oplus(\mathbb Z/81)^3.
\end{split}
\tag{A.22.3}
\]
Every nonempty lifting fiber from class two to class three is
a torsor under (A.22.3).
\end{statement}
\begin{proof}
We verify the integral lattice in (A.22.1). In a metabelian
graded Lie ring the adjoint generators commute on the derived
ideal. Antisymmetry and Jacobi collect every degree-$r$ bracket
into left-normed brackets
\[
 [x_i,x_j,x_{k_3},\ldots,x_{k_r}],\qquad
 i>j,\quad j\le k_3\le\cdots\le k_r.                 \tag{A.22.4}
\]
To prove independence integrally, use
$\Lambda=\mathbb Z_3[[X_1,\ldots,X_{10}]]$ and the upper
triangular matrices with diagonal $(1+X_i,1)$ and upper-right
entry $e_i\in\Lambda^{10}$. Their closed generated group is
pro-three in the $(3,X_1,\ldots,X_{10})$-adic topology, so
the free pro-three universal property defines a continuous
homomorphism from $N$. Commutators are translations, so it
kills $N''$. The translation of $[x_i,x_j]$ is
$X_i e_j-X_j e_i$; each further right bracket multiplies
its leading term by $-X_k$.

In a fixed multidegree of total degree $r$, let $J$ be the
variable support. The multiplication kernel consists of
coefficient vectors indexed by $J$ with sum zero. The brackets
(A.22.4) in that multidegree map, up to their common sign, to
the standard differences with distinguished index $j=\min J$
and other index $i\in J\setminus\{j\}$. They are an integral
basis. Collection proves spanning and these matrices prove
independence, hence the absence of three-torsion and the
asserted isomorphism. The bracket/multiplication construction
is natural in $V$, proving Galois equivariance. Subtracting
the symmetric-power ranks gives the displayed rank formula.

Over $K_0$ these finite coefficient representations have
three-group image. The composition-series proof of A.1 gives
$H^i(G_S,M_r/3^n)=0$ for $i\ge2$. Normalized cocycles of
A.10 give the finite-$S$ complex
$[M_r^c\xrightarrow{u-1}M_r]$. For $r\ge3$, the weights
give injective $D_{r,q}$, as in F.5, so the integral
degree-one residue invariants vanish. Localization and
F.1--F.3 prove (A.22.2), including the nonzero derived limit.

The metabelian relation starts in degree four. The marked
relation also does: modulo $\gamma_3N$, conjugates of the
marked cusp generator differ from it by central elements,
so its normal closure is abelian there. Thus
$[I_{\rm mark},I_{\rm mark}]\subset\gamma_3N$ and
$\gamma_3(I_{\rm mark})\subset\gamma_4N$.
This proves the common degree-three coefficient.

The exact program \texttt{code/a\_degree3\_lattice.py}
constructs the $330$-dimensional multiplication kernel by
these multidegree differences, verifies the induced matrices,
and gives involution eigenspace ranks $164,166$.
The $330$ by $164$ matrix $(u-1)|M_3^c$ has Smith valuations
\[
 0^{73},\quad1^7,\quad2^{56},\quad3^{25},\quad4^3.
\]
All 164 pivots are obtained modulo $3^{10}$; their valuations
are at most four. The recorded invertible row and column
operations therefore certify the exact three-adic Smith
factors, giving (A.22.3). C.1 makes the target geometric
centralizer trivial, so no additional stabilizer identifies
the central lifting torsor.
\end{proof}

\begin{statement}{A.23}{Proved}{The complete higher images and fibers in Task A.6}
For each of $G=N/N''$ and $G=N/\gamma_3(I_{\rm mark})$,
let $X_n(s,G)$, $G_n$, $A_n$ and $J_n$ be the explicitly
constructed finite arithmetic data of A.43 for an exact
class-two target $s$. Then the image consists exactly of
the classes satisfying
\[
             X_n(s,G)\ne\varnothing\quad\text{for every }n,
\]
and its entire fiber is
\[
                  \varprojlim_n X_n(s,G)/J_n^c.
\]
Every finite test terminates and uses the actual based
arithmetic action and finite target table. The inherited
ramification constraints give its transition maps, including
for infinitely ramified targets. The lower-central quotients
of A.21 are cofinal for the same section problem. Their
metabelian coefficient lattices and derived limits are A.22;
their marked coefficient lattices and derived limits are A.44.
The exact maps to the inherited quotients, with both kernels,
are (A.44.2).
\end{statement}
\begin{proof}
The two geometric kernels are closed and arithmetic-stable.
The quotient maps in (A.43.2) are their actual finite images,
and (A.43.8) identifies their inverse limits with the specified
groups. Normalization at complex conjugation is possible
inside each pro-three lifting kernel. A.43 proves the
uniform ramification bound from the abelian target and
enumerates every possible finite graph field within its
explicit degree and discriminant bounds. The resulting
finite sets therefore give both directions of the image
criterion, with continuous compatible cocycles.

For the fiber, A.43 proves that the image of the actual
geometric kernel is $J_n$, and that its $c$-fixed subgroup
surjects onto $J_n^c$. The global class-two centralizer is
trivial. Its compact conjugator comparison consequently
gives precisely the displayed inverse limit of orbits.
No extra finite-level centralizer is divided out.

A.21 proves both the group and section-class comparisons
with the stated finite-rank lower-central towers. A.22 and
A.44 identify the actual coefficient lattices and actions;
their residue formulas use the image on restricted integral
families and retain its nonzero first derived limit at every
positive higher layer. Thus convergence has not been
replaced by independent local solvability. Formula (A.44.2)
comes from the same based Fermat covering and gives the
requested inherited comparisons. These statements together
settle every image, fiber, cofinality and coefficient-layer
requirement of Task A.6. The finite arithmetic stages form
an infinite criterion; no total stopping algorithm is used.
\end{proof}

\subsection{Explicit nine-torsion covers at the first auxiliary Frobenius}

\begin{statement}{A.36}{Proved}{Miller covers and their two exact norm comparisons at seven}
The Fermat curve in the task has the following specified isomorphism
with the minimal Weierstrass curve $E_m:y_m^2+y_m=x_m^3-7$:
\[
 x_m=\frac{-3z}{x+y},\quad
 y_m=\frac{4x-5y}{x+y},\qquad
 [x:y:z]=[y_m+5:4-y_m:-3x_m].                 \tag{A.36.1}
\]
It carries $O$ to the Weierstrass origin and carries
$P,Q_c$ to $(3,-5),(3,4)$, respectively.
All of $E_m[9]$ is defined over the unramified degree-three
extension of $\Q_7$. Its reduction over
$k=\F_7[\theta]/(\theta^3-2)$ has exactly 81 points.
One point above $P$ under multiplication by three is
\[
 Q=(2+\theta^2,\ 1+\theta+6\theta^2),\qquad
 \varphi Q-Q=(6,2),                           \tag{A.36.2}
\]
where $\varphi$ raises the residue coordinates to their seventh
powers. The unique nine-torsion lifts give the corresponding
identities over the unramified local extension.

For each of the nine points $Q$ above either $d=P$ or $d=Q_c$,
let $v_R=x_m-x_m(R)$ and let $\ell_{R,S}$ be the line through
$R,S$, tangent at $R$ when $R=S$, normalized to have
$y_m$ coefficient one. Then the actual function
\[
 a_Q=\frac{\ell_{Q,Q}\ell_{2Q,Q}}{v_{2Q}v_d},\qquad
 \operatorname{div}(a_Q)=3(Q)-(d)-2(O),        \tag{A.36.3}
\]
defines a connected cubic cover $w^3=a_Q$ which is finite
\'etale over $U$. Its geometric character is an explicit
Kummer class with nonzero cusp valuations at $d,O$.
With $\iota(R)=-R$, the following function identities hold:
\[
 a_Q\,\iota^*a_Q=\frac{(x_m-x_m(Q))^3}{x_m-x_m(d)},
 \qquad
 \prod_{i=0}^2 a_{\varphi^iQ}=g_Q^3.           \tag{A.36.4}
\]
In the second identity, put $Q_i=\varphi^iQ$,
$S=Q_0+Q_1$, and
\[
       g_Q=\frac{\ell_{Q_0,Q_1}\ell_{S,Q_2}}{v_Sv_d}.
                                                        \tag{A.36.5}
\]
This function is defined over $\Q_7$, and its reduction is
defined over $\F_7$.

For the exact Fermat parameter $s=z/x$ and its unit tangent,
the regularized constants are
\[
 (s^2a_Q)(O)=9,
 \qquad (s^2(x_m-3))(O)=(s^2(x_m-3\zeta^2))(O)=9.
                                                        \tag{A.36.6}
\]
In the large model $Y^2=X^3-432$, $X=4x_m$, $Y=8y_m+4$,
the corresponding normalized Miller function is $4a_Q$,
with regularized constant 36. For the residue embedding
$\zeta\mapsto2\in\F_7$, the cubic residue of 9 is
$4=\zeta^2$, whereas that of 36 is 1.
These are arithmetic cover data; no central Galois action
on an intrinsic Heisenberg quotient is inferred from them.
\end{statement}
\begin{proof}
Substitution in (A.36.1) gives the identity
\[
 (x+y)^3\bigl(y_m^2+y_m-x_m^3+7\bigr)
                     =27(x^3+y^3+z^3).
\]
The displayed inverse gives the reverse identity.
The rational maps extend across their omitted points because
both projective curves are smooth. Their composites on the
displayed dense affine opens are the identity, proving the
claimed isomorphism and its stated cusp images.
The discriminant of the minimal equation is $-3^9$,
so it has good reduction at seven. Multiplication by nine
is \'etale over $\Z_7$, and reduction identifies its
nine-torsion over an unramified extension with the
nine-torsion of the residue curve.

The exact program \texttt{a\_miller\_f343\_root.py}
constructs $k$, lists all points of $y_m^2+y_m=x_m^3$,
and applies the stated Weierstrass group law. There are
324 points, 81 killed by nine, and nine killed by three.
Every three-torsion point has coordinates in $\F_7$.
It lists all nine preimages of each specified cusp and
verifies (A.36.2) and each Frobenius orbit.
The geometric nine-torsion has order 81 because nine
is invertible in characteristic seven; the enumerated
points therefore exhaust it. Finite \'etale lifting
proves the local assertion, rather than merely an
assertion about an unrelated finite-field group.

The divisors of the normalized line and vertical functions
are
\[
 \begin{split}
 \operatorname{div}(\ell_{R,S})
       &=(R)+(S)+(-R-S)-3(O),\\
 \operatorname{div}(v_R)&=(R)+(-R)-2(O).
 \end{split}
\]
Their substitution proves (A.36.3). At every point of $U$
the valuation is a multiple of three, so the normalized
cubic cover is unramified there. Characteristic zero
then makes it \'etale. Its valuation $-1$ at $d$ proves
that the function is not a cube and that the cover is
connected. This gives an actual finite cover whose Kummer character
is defined by the displayed function.

Writing $t=-x_m/y_m$ at the origin gives
$x_m=t^{-2}+O(t)$, $y_m=-t^{-3}+O(1)$.
Each normalized line has leading term $-t^{-3}$ and
each vertical function has leading term $t^{-2}$.
Consequently $a_Q$ and $g_Q$ have leading term $t^{-2}$.
The divisors of the two sides of the first identity
in (A.36.4) agree, and their leading coefficients at
$O$ are both one. Their ratio is therefore the constant
one on the smooth proper curve, proving that identity.

For the second identity, write
$D_Q=\varphi Q-Q$. The finite computation and finite
\'etale comparison give $D_Q\in E_m[3](\Q_7)$.
Thus $Q_i=Q+iD_Q$, and
$Q_0+Q_1+Q_2=3Q+3D_Q=d$.
The line-divisor calculation in (A.36.5) gives
\[
       \operatorname{div}(g_Q)
                  =(Q_0)+(Q_1)+(Q_2)-(d)-2(O).
\]
This divisor is fixed by $\varphi$. Its normalized
leading coefficient is one, so $\varphi(g_Q)=g_Q$.
The function descends to $\Q_7$. Comparing divisors
and the same leading coefficient proves the cube
identity. The finite-field program also multiplies
the function numerators in the exact ring
$k[x_m,y_m]/(y_m^2+y_m-x_m^3)$ and verifies both
identities by cross-multiplication, for all 18 points.

Finally on the Fermat chart $x=1$,
$y=-(1+s^3)^{1/3}$ gives
$x_m=9s^{-2}+O(s)$ and $y_m=-27s^{-3}+O(1)$.
Hence $s=3t+O(t^4)$, which proves (A.36.6).
The large-model line and vertical functions are eight
and four times the corresponding minimal-model
functions, so their two-line quotient is $4a_Q$.
Reduction of the retained constants gives
$9\equiv2$, $36\equiv1\pmod7$; the cubic residue
map on $\F_7^\times$ is $u\mapsto u^2$.
This proves the stated tangent phases without dropping
the origin term. The program retains the characteristic-zero
coordinate identity, tangent expansions, complete field
coordinates, line coefficients, and all exact outputs.
\end{proof}

\begin{statement}{A.37}{Proved}{An order-eighty-one arithmetic quotient and a tame-letter Frobenius coefficient at seventeen}
Put $K_0=\mathbb Q(\zeta)$, $\zeta^2+\zeta+1=0$, and
$L_{17}=K_0(\beta)$ with $\beta^3=17$. Let
$\sigma(\beta)=\zeta\beta$ and set
\[
 \delta=\frac{-7-\zeta-3\beta-\beta^2}{3},\qquad
 \delta_i=\sigma^i(\delta),\qquad
 r_i^3=\delta_i\quad(i\in\mathbb Z/3).
 \tag{A.37.1}
\]
The extension
\[
 M_{17}=L_{17}(r_0,r_1,r_2)/K_0
 \tag{A.37.2}
\]
is Galois, unramified outside three and seventeen, with group
\[
 G=(\mathbb F_3)^3\rtimes\langle s\rangle,
 \qquad s^3=1,
 \tag{A.37.3}
\]
where $s(r_i)=r_{i+1}$ and $s(\beta)=\zeta\beta$.
For $a=(a_0,a_1,a_2)$ write
$\tau_a(r_i)=\zeta^{a_i}r_i$, $\tau_a(\beta)=\beta$.
Then
$s\tau_as^{-1}=\tau_{Sa}$, where $(Sa)_i=a_{i-1}$.
The group has order $81$, lower-central orders
$81,9,3,1$, exponent nine, and trivial fourth
three-Zassenhaus subgroup.

Use the arithmetic generators $m,n,t$ of A.28 and A.31,
with the exact Kummer normalization of A.35. Their
modulo-three characters are
\[
 \begin{array}{c|ccc}
       &m&n&t\\\hline
 \kappa_\zeta&1&0&0\\
 \kappa_{17}&0&0&1
 \end{array}.
 \tag{A.37.4}
\]
The restriction map to (A.37.2), composed with the
actual free pro-three presentation, is a continuous
epimorphism
\[
 F_3(m,n,t)\longrightarrow
 G_{K_0,\{3,17\}}(3)\longrightarrow G.
 \tag{A.37.5}
\]
Let $W\in D_3(F_3(m,n,t))$ be the actual arithmetic
Frobenius word of A.30. Under (A.37.5),
\[
 W\longmapsto\tau_{(2,2,2)},\qquad
 [t,[t,m]]\longmapsto\tau_{(1,1,1)}.
 \tag{A.37.6}
\]
In the restricted degree-three Lie algebra
$D_3(F_3)/D_4(F_3)$, this map annihilates the three
generator cubes and every ordinary degree-three
component except the multidegree $(1,0,2)$ component
generated by $[t,[t,m]]$. Consequently the coefficient
of this bracket in the full $W$ is exactly $2\in\mathbb F_3$.
This records a coefficient involving tame inertia;
the map killing $t$ in A.31 does not retain it.
\par\smallskip\noindent\emph{Manuscript comparison:
\S51.16, folios 733v--736v, scans 172--178,
formulas (371)--(379); Tasks A.6 and B.4.}
\end{statement}
\begin{proof}
For $a,b,c\in K_0$ direct multiplication of the three
conjugates gives
\[
 N_{L_{17}/K_0}(a+b\beta+c\beta^2)
       =a^3+17b^3+289c^3-51abc.
 \tag{A.37.7}
\]
At $(a,b,c)=(-7-\zeta,-3,-1)$ the right side is
$27\zeta$. Thus
\[
                    \delta_0\delta_1\delta_2=\zeta.
 \tag{A.37.8}
\]
Each $\delta_i$ is integral away from three. Since
its norm is a unit, it is a unit at every finite
prime away from three. Kummer extensions by its
cube root are unramified at such primes, because
three is invertible there and the radical is a unit.
The extension $L_{17}/K_0$ is ramified only at three
and seventeen. This proves the asserted ramification
bound once the degree and Galois group are established.

The three classes of $\delta_i$ are independent in
$L_{17}^{\times}/L_{17}^{\times3}$. Indeed the map
\[
 \mathbb F_3[\langle\sigma\rangle]
      \longrightarrow L_{17}^{\times}/L_{17}^{\times3},
 \qquad \sigma^i\longmapsto[\delta_i]
 \tag{A.37.9}
\]
is a module homomorphism. Every nonzero submodule
of the regular module
$\mathbb F_3[u]/(u-1)^3$ contains its socle,
generated by $1+u+u^2$. By (A.37.8) that element
maps to $[\zeta]$. The latter is nonzero: adjoining
a cube root of $\zeta$ gives $K_0(\mu_9)$, a cubic
extension unramified at seventeen, whereas
$L_{17}/K_0$ is Eisenstein at the unique prime over
seventeen. These two cubic extensions therefore have
different ramification indices and cannot coincide.
Thus the kernel in (A.37.9) is zero.

Kummer theory now gives
$\operatorname{Gal}(M_{17}/L_{17})=(\mathbb F_3)^3$
with the displayed diagonal actions. Independence
also proves that the defining radical algebra is a
field of degree $27$ over $L_{17}$. Cyclic permutation
of its generators defines the specified extension
$s$ of $\sigma$, with $s^3=1$. This proves the
semidirect presentation (A.37.3). Its commutators are
\[
 \gamma_2G=(S-1)(\mathbb F_3)^3,\quad
 \gamma_3G=(S-1)^2(\mathbb F_3)^3
       =\langle\tau_{(1,1,1)}\rangle,\quad
 \gamma_4G=1.
 \tag{A.37.10}
\]
The three dimensions are $2,1,0$. Every cube in
$G$ is diagonal central, while $G^9=1$ and
$(\gamma_2G)^3=1$. The defining formula
$D_j(G)=\prod_{i3^a\ge j}(\gamma_iG)^{3^a}$
therefore gives $D_4(G)=1$.

We next compute the actual local Frobenius. The
residue field of $K_0$ at seventeen is
$k=\mathbb F_{289}=\mathbb F_{17}[\zeta]$.
At its unique extension to $L_{17}$ the radical
$\beta$ reduces to zero, so every $\delta_i$ has
the same residue
\[
 d_0=(-7-\zeta)/3=9+11\zeta,\qquad
 d_0^3=\zeta,\qquad d_0^{96}=\zeta^2.
 \tag{A.37.11}
\]
In particular $d_0$ is not a cube in $k$.
The polynomial $R^3-d_0$ defines the unramified
residue extension of degree three. Its roots obey
\[
                       R^{289}=d_0^{96}R=\zeta^2R.
 \tag{A.37.12}
\]
An arithmetic Frobenius fixing the chosen uniformizer
roots fixes $\beta$ and acts by $r_i\mapsto\zeta^2r_i$
for all three radicals. To see that the residue
calculation determines this exact automorphism,
the ratio of any two possible images of $r_i$ is
in $\mu_3$, whose reduction is injective at seventeen.
This proves the first formula in (A.37.6).
It is independent of conjugating the chosen
decomposition subgroup because this image is central.

Finally put $v=r_0r_1r_2$, so $v^3=\zeta$.
For an element $\tau_as^k$, the two characters
$\kappa_\zeta,\kappa_{17}$ are respectively
$a_0+a_1+a_2$ and $k$. Thus the actual images of
the normalized generators have the form
\[
 m\mapsto\tau_a,\qquad n\mapsto\tau_b,
 \qquad t\mapsto\tau_cs,
 \qquad \sum a_i=1,\quad\sum b_i=\sum c_i=0.
 \tag{A.37.13}
\]
We have made no choice of their individual coordinates.
All three images have cube one. Conjugation by
$\tau_cs$ on the abelian kernel is $S$, and
\[
 (S-1)^2a=(1+S+S^2)a=(1,1,1),\qquad
                         (S-1)^2b=0.
 \tag{A.37.14}
\]
Every bracket with two entries from the abelian
kernel vanishes. The only surviving ordinary
degree-three bracket is therefore $[t,[t,m]]$;
its image is the second formula in (A.37.6).
The standard free restricted degree-three basis
consists of the ordinary degree-three free Lie
part and the three generator cubes. The preceding
calculations identify the entire map on that basis,
so the Frobenius image proves that its indicated
coefficient is two. This argument uses the full
three-generator word and the actual arithmetic
restriction (A.37.5).

For comparison with the lower Heisenberg test, the
fixed field of the diagonal subgroup in (A.37.10)
is
\[
 L_{17}(v,w),\qquad
 v^3=\zeta,\quad w=r_0^2r_1,\quad
 w^3=\delta_0^2\delta_1.
 \tag{A.37.15}
\]
The two character rows $(1,1,1)$ and $(2,1,0)$
are independent, so this field has degree nine
over $L_{17}$, hence degree $27$ over $K_0$.
The action satisfies $s(w)=(v/\delta_0)w$.
Its weighted radical has residue $d_0^3=\zeta$,
which already has the cube root $d_0$ in $k$.
The Frobenius image computed above becomes the
identity in this quotient, exactly as follows
from its diagonal kernel. This specifies both
the lower quotient map and the coefficient it
forgets.

The exact program \texttt{a\_wreath17\_root.py}
verifies (A.37.7), the residue field and its actual
degree-three radical extension, all $81$ group
elements, the lower-central orders, and every one
of the $729$ triples in (A.37.13). The preliminary
norm search and its exact algebraic certification
are retained in \texttt{a\_norm\_zeta\_17.py}.
\end{proof}
